\chapter{2005年湖南卷（文）}
\section{单选题}
\begin{problem}
设全集 \(U = \{-2, -1, 0, 1, 2\}\)，\(A = \{-2, -1, 0\}\)，\(B = \{0, 1, 2\}\)，则 \((\complement_U A) \cap B =\)（\quad）

\choices
    {\( \{0\} \)}
    {\(\{-2, -1\}\)}
    {\(\{1, 2\}\)}
    {\(\{0, 1, 2\}\)}
\end{problem}

\begin{answer}
C
\end{answer}

\begin{problem}
\(\tan 600^{\circ}\) 的值是（\quad）

\choices
    {\(-\frac{\sqrt{3}}{3}\)}
    {\(\frac{\sqrt{3}}{3}\)}
    {\(-\sqrt{3}\)}
    {\(\sqrt{3}\)}
\end{problem}

\begin{answer}
D
\end{answer}

\begin{problem}
函数 \( f(x) = \sqrt{1 - 2^x} \) 的定义域是（\quad）

\choices
    {\((- \infty, 0]\)}
    {\([0, +\infty)\)}
    {\((-\infty, 0)\)}
    {\((-\infty, +\infty)\)}
\end{problem}

\begin{answer}
A
\end{answer}

\begin{problem}
如图，正方体 \(ABCD - A_{1}B_{1}C_{1}D_{1}\) 的棱长为 \(1\)，\(E\) 是 \(A_{1}B_{1}\) 的中点. 则 \(E\) 到平面 \(ABC_{1}D_{1}\) 的距离为（\quad）

\bitmapfigure[max width=0.42\linewidth]{img/2005/hunan_liberal/q4_fig1.png}

\choices
    {\(\frac{\sqrt{3}}{2}\)}
    {\(\frac{\sqrt{2}}{2}\)}
    {\(\frac{1}{2}\)}
    {\(\frac{\sqrt{3}}{3}\)}
\end{problem}

\begin{answer}
B
\end{answer}

\begin{problem}
已知数列 \(\{a_{n}\}\) 满足 \(a_{1} = 0\)，\(a_{n + 1} = \frac{a_{n} - \sqrt{3}}{\sqrt{3}a_{n} + 1}\)（\(n \in \N^{*}\)），则 \(a_{20} =\)（\quad）

\choices
    {0}
    {\(-\sqrt{3}\)}
    {\(\sqrt{3}\)}
    {\(\frac{\sqrt{3}}{2}\)}
\end{problem}

\begin{answer}
B
\end{answer}

\begin{problem}
设集合 \(A = \left\{x \mid \frac{x - 1}{x + 1} < 0\right\}\)，\(B = \{x \mid |x - 1| < a\}\)，则“\(a = 1\)”是“\(A \cap B \neq \varnothing\)”的（\quad）

\choices
    {充分不必要条件}
    {必要不充分条件}
    {充要条件}
    {既不充分又不必要条件}
\end{problem}

\begin{answer}
A
\end{answer}

\begin{problem}
设直线的方程是 \(Ax + By = 0\)，从 \(1\)，\(2\)，\(3\)，\(4\)，\(5\) 这五个数中每次取两个不同的数作为 \(A\)，\(B\) 的值，则所得不同直线的条数是（\quad）

\choices
    {20}
    {19}
    {18}
    {16}
\end{problem}

\begin{answer}
C
\end{answer}

\begin{problem}
已知双曲线 \(\frac{x^2}{a^2} - \frac{y^2}{b^2} = 1\)（\(a > 0\)，\(b > 0\)）的右焦点为 \(F\)，右准线与一条渐近线交于点 \(A\)，\(\triangle OAF\) 的面积为 \(\frac{a^2}{2}\)（\(O\) 为原点），则两条渐近线的夹角为（\quad）

\choices
    {\(30^{\circ}\)}
    {\(45^{\circ}\)}
    {\(60^{\circ}\)}
    {\(90^{\circ}\)}
\end{problem}

\begin{answer}
D
\end{answer}

\begin{problem}
\(P\) 是 \(\triangle ABC\) 所在平面上一点，若 \(\overrightarrow{PA} \cdot \overrightarrow{PB} = \overrightarrow{PB} \cdot \overrightarrow{PC} = \overrightarrow{PC} \cdot \overrightarrow{PA}\)，则 \(P\) 是 \(\triangle ABC\) 的（\quad）

\choices
    {外心}
    {内心}
    {重心}
    {垂心}
\end{problem}

\begin{answer}
D
\end{answer}

\begin{problem}
某公司在甲、乙两地销售一种品牌车，利润（单位：万元）分别为 \(L_{1} = 5.06x - 0.15x^{2}\) 和 \(L_{2} = 2x\)，其中 \(x\) 为销售量（单位：辆）. 若该公司在这两地共销售 \(15\) 辆车，则能获得的最大利润为（\quad）

\choices
    {45.606}
    {45.6}
    {45.56}
    {45.51}
\end{problem}

\begin{answer}
B
\end{answer}

\section{填空题}
\begin{problem}
设直线 \(2x + 3y + 1 = 0\) 和圆 \(x^{2} + y^{2} - 2x - 3 = 0\) 相交于点 \(A\)，\(B\)，则弦 \(AB\) 的垂直平分线方程是\fillinblank{}.
\end{problem}

\begin{answer}
\(3x - 2y - 3 = 0\)
\end{answer}

\begin{problem}
一工厂生产了某种产品 \(16800\) 件，它们来自甲、乙、丙 \(3\) 条生产线，为检查这批产品的质量，决定采用分层抽样的方法进行抽样. 已知甲、乙、丙三条生产线抽取的个体数组成一个等差数列，则乙生产线生产了\fillinblank{}件产品.
\end{problem}

\begin{answer}
5600
\end{answer}

\begin{problem}
在 \((1 + x) + (1 + x)^2 + \cdots + (1 + x)^6\) 的展开式中，\(x^2\) 项的系数是\fillinblank{}. （用数字作答）
\end{problem}

\begin{answer}
35
\end{answer}

\begin{problem}
设函数 \( f(x) \) 的图象关于点 \((1,2)\) 对称，且存在反函数 \( f^{-1}(x) \)，\( f(4) = 0 \)，则 \( f^{-1}(4) = \)\fillinblank{}.
\end{problem}

\begin{answer}
\(-2\)
\end{answer}

\begin{problem}
已知平面 \(\alpha\)，\(\beta\) 和直线 \(m\)，给出条件：

\begin{circlelist}
\item \( m \parallel \alpha \)；
\item \( m \perp \alpha \)；
\item \( m \subset \alpha \)；
\item \( \alpha \perp \beta \)；
\item \( \alpha \parallel \beta \).
\end{circlelist}

\begin{enumerate}
\item 当满足条件\fillinblank{} 时，有 \( m \parallel \beta \)；
\item 当满足条件\fillinblank{} 时，有 \( m \perp \beta \). （填所选条件的序号）
\end{enumerate}
\end{problem}

\begin{answer}
\(\circled{3}\circled{5}\)；\(\circled{2}\circled{5}\)
\end{answer}

\section{解答题}
\begin{problem}
已知数列 \(\{\log_2(a_n - 1)\}\)（\(n \in \N^*\)）为等差数列，且 \(a_1 = 3\)，\(a_3 = 9\).

\begin{enumerate}
\item 求数列 \( \{a_{n}\} \) 的通项公式；
\item 证明： \( \frac{1}{a_{2}-a_{1}}+\frac{1}{a_{3}-a_{2}}+\cdots+\frac{1}{a_{n+1}-a_{n}}<1 \).
\end{enumerate}

\begin{solution}
\begin{enumerate}
\item
令 \(d_n=\log_2(a_n-1)\). 由题意，\(\{d_n\}\) 是等差数列. 因为
\(d_1=\log_2(3-1)=1\)，\(d_3=\log_2(9-1)=3\)，所以公差为
\(\frac{d_3-d_1}{2}=1\)，从而 \(d_n=1+(n-1)=n\). 因此
\(a_n-1=2^{d_n}=2^n\)，即
\[
a_n=2^n+1.
\]

\item
由上式，对 \(k\in\N^*\) 有
\(a_{k+1}-a_k=2^{k+1}+1-(2^k+1)=2^k\). 所以
\[
\frac{1}{a_2-a_1}+\frac{1}{a_3-a_2}+\cdots+\frac{1}{a_{n+1}-a_n}
=\frac12+\frac14+\cdots+\frac{1}{2^n}
=1-\frac{1}{2^n}<1.
\]
故不等式成立.
\end{enumerate}
\end{solution}
\end{problem}

\begin{problem}
已知在 \(\triangle ABC\) 中，\(\sin A(\sin B + \cos B) - \sin C = 0\)，\(\sin B + \cos 2C = 0\)，求角 \(A\)，\(B\)，\(C\) 的大小.

\begin{solution}
因为 \(A\)，\(B\)，\(C\) 是三角形的内角，所以
\(0<A,B,C<\pi\)，且 \(A+B+C=\pi\). 由三角形内角和，
\[
\sin C=\sin(A+B)=\sin A\cos B+\cos A\sin B.
\]
代入第一个已知等式，得
\[
\sin A\sin B+\sin A\cos B
=\sin A\cos B+\cos A\sin B,
\]
即 \(\sin B(\sin A-\cos A)=0\). 由于 \(0<B<\pi\)，\(\sin B>0\)，所以
\(\sin A=\cos A\)，从而 \(A=\frac{\pi}{4}\).

于是 \(B+C=\frac{3\pi}{4}\). 由第二个已知等式，
\[
\sin\left(\frac{3\pi}{4}-C\right)=-\cos 2C
=\sin\left(2C-\frac{\pi}{2}\right).
\]
因此存在 \(k\in\Z\)，使得
\[
\begin{cases}
\dfrac{3\pi}{4}-C=2C-\dfrac{\pi}{2}+2k\pi,\\
\dfrac{3\pi}{4}-C=\pi-\left(2C-\dfrac{\pi}{2}\right)+2k\pi.
\end{cases}
\]
第二种情形给出 \(C=\frac{3\pi}{4}+2k\pi\)，不满足
\(0<C<\frac{3\pi}{4}\). 第一种情形给出
\(C=\frac{5\pi}{12}-\frac{2k\pi}{3}\)，结合
\(0<C<\frac{3\pi}{4}\) 只有 \(k=0\)，故 \(C=\frac{5\pi}{12}\). 于是
\[
B=\frac{3\pi}{4}-\frac{5\pi}{12}=\frac{\pi}{3}.
\]
所以 \(A=45^{\circ}\)，\(B=60^{\circ}\)，\(C=75^{\circ}\).
\end{solution}
\end{problem}

\begin{problem}
如图 1，已知 \(ABCD\) 是上、下底边长分别为 \(2\) 和 \(6\)，高为 \(\sqrt{3}\) 的等腰梯形，将它沿对称轴 \(OO_{1}\) 折成直二面角，如图 2.

\begin{enumerate}
\item 证明： \(AC \perp BO_{1}\)；
\item 求二面角 \(O - AC - O_{1}\) 的大小.
\end{enumerate}

\bitmapfigure[max width=0.46\linewidth]{img/2005/hunan_liberal/q18_fig1.png}

\begin{solution}
折叠后建立直角坐标系，取折痕 \(OO_1\) 为 \(y\) 轴，使
\[
O(0,0,0),\quad O_1(0,\sqrt{3},0),\quad A(-3,0,0),\quad D(-1,\sqrt{3},0),
\]
\[
B(0,0,3),\quad C(0,\sqrt{3},1).
\]
其中左、右两半分别位于两个互相垂直的平面内，且这些坐标保持了折叠前各边的长度.

\begin{enumerate}
\item 由坐标得
\[
\overrightarrow{AC}=(3,\sqrt{3},1),\qquad
\overrightarrow{BO_1}=(0,\sqrt{3},-3).
\]
因此
\[
\overrightarrow{AC}\cdot\overrightarrow{BO_1}=3-3=0,
\]
故 \(AC\perp BO_1\).
\item 平面 \(OAC\) 和平面 \(O_1AC\) 的一个法向量分别为
\[
\bs{n}_1=\overrightarrow{AO}\times\overrightarrow{AC}=(0,-3,3\sqrt{3}),
\qquad
\bs{n}_2=\overrightarrow{AO_1}\times\overrightarrow{AC}=(\sqrt{3},-3,0).
\]
设二面角 \(O-AC-O_1\) 的大小为 \(\theta\)，则
\[
\cos\theta=\frac{|\bs{n}_1\cdot\bs{n}_2|}{|\bs{n}_1||\bs{n}_2|}
=\frac{9}{6\cdot2\sqrt{3}}=\frac{\sqrt{3}}{4}.
\]
所以二面角 \(O-AC-O_1\) 的大小为
\(\arccos\frac{\sqrt{3}}{4}\).
\end{enumerate}
\end{solution}
\end{problem}

\begin{problem}
设 \(t \neq 0\)，点 \(P(t,0)\) 是函数 \(f(x) = x^3 + ax\) 与 \(g(x) = bx^2 + c\) 的图象的一个公共点，两函数的图象在点 \(P\) 处有相同的切线.

\begin{enumerate}
\item 用 \(t\) 表示 \(a\)，\(b\)，\(c\)；
\item 若函数 \( y = f(x) - g(x) \) 在 \((-1, 3)\) 上单调递减，求 \( t \) 的取值范围.
\end{enumerate}

\begin{solution}
\begin{enumerate}
\item 记 \(h(x)=f(x)-g(x)\). 因为 \(P(t,0)\) 在两条曲线上，所以
\[
t^3+at=0,\qquad bt^2+c=0.
\]
又因为两条曲线在 \(P\) 处有相同的切线，所以切线斜率相等，即
\[
f'(t)=g'(t),\qquad 3t^2+a=2bt.
\]
由 \(t\neq0\)，第一式给出 \(a=-t^2\). 将其代入斜率关系，得
\(2t^2=2bt\)，所以 \(b=t\). 再由 \(bt^2+c=0\)，得 \(c=-t^3\). 因此
\[
a=-t^2,\qquad b=t,\qquad c=-t^3.
\]

\item 代入所得参数，
\[
h(x)=x^3-tx^2-t^2x+t^3=(x-t)^2(x+t),
\]
从而
\[
h'(x)=3x^2-2tx-t^2=(3x+t)(x-t).
\]
函数 \(h(x)\) 在 \((-1,3)\) 上单调递减，要求 \(h'(x)\leqslant0\) 在该区间上成立. 因为 \(h'(x)\) 是开口向上的二次函数，所以区间 \((-1,3)\) 必须位于两个零点之间.

当 \(t>0\) 时，两个零点为 \(-\frac{t}{3}\) 和 \(t\)，且 \(-\frac{t}{3}<t\). 要使 \((-1,3)\) 位于两根之间，需
\[
-\frac{t}{3}\leqslant-1,\qquad t\geqslant3,
\]
这等价于 \(t\geqslant3\).

当 \(t<0\) 时，两个零点为 \(t\) 和 \(-\frac{t}{3}\)，且 \(t<-\frac{t}{3}\). 此时需
\[
t\leqslant-1,\qquad -\frac{t}{3}\geqslant3,
\]
这等价于 \(t\leqslant-9\).

当 \(t=3\) 或 \(t=-9\) 时，导数为零的点分别是区间端点，开区间 \((-1,3)\) 内仍有 \(h'(x)<0\). 故 \(t\) 的取值范围为
\[
(-\infty,-9]\cup[3,+\infty).
\]
\end{enumerate}
\end{solution}
\end{problem}

\begin{problem}
某单位组织 \(4\) 个部门的职工旅游，规定每个部门只能在韶山、衡山、张家界 \(3\) 个景区中任选一个，假设各部门选择每个景区是等可能的.

\begin{enumerate}
\item 求 \(3\) 个景区都有部门选择的概率；
\item 求恰有 \(2\) 个景区有部门选择的概率.
\end{enumerate}

\begin{solution}
每个部门有 \(3\) 种选择，且各部门的选择相互独立，所以所有选择的总数为
\(3^4\).

\begin{enumerate}
\item 若 \(3\) 个景区都有部门选择，则一个景区有 \(2\) 个部门选择，另外两个景区各有 \(1\) 个部门选择. 先选出有 \(2\) 个部门的景区，有 \(3\) 种；再从 \(4\) 个部门中选出这 \(2\) 个部门，有 \(\mathrm C_4^2\) 种；剩下的两个部门分别选择另外两个景区，有 \(2\) 种. 因而所求概率为
\[
\frac{3\mathrm C_4^2\cdot2}{3^4}=\frac{36}{81}=\frac{4}{9}.
\]
\item 若恰有 \(2\) 个景区有部门选择，先选出这 \(2\) 个景区，有 \(\mathrm C_3^2\) 种. 对固定的两个景区，所有部门在其中选择共有 \(2^4\) 种，其中全部选择第一个或全部选择第二个的情形各有 \(1\) 种，故恰好使用两个景区有 \(2^4-2\) 种. 因而所求概率为
\[
\frac{\mathrm C_3^2(2^4-2)}{3^4}=\frac{42}{81}=\frac{14}{27}.
\]
\end{enumerate}
\end{solution}
\end{problem}

\begin{problem}
已知椭圆 \(C: \frac{x^2}{a^2} + \frac{y^2}{b^2} = 1\)（\(a > b > 0\)）的左、右焦点为 \(F_1\)，\(F_2\)，离心率为 \(e\). 直线 \(l: y = ex + a\) 与 \(x\) 轴，\(y\) 轴分别交于点 \(A\)，\(B\)，\(M\) 是直线 \(l\) 与椭圆 \(C\) 的一个公共点，\(P\) 是点 \(F_1\) 关于直线 \(l\) 的对称点，设 \(\overrightarrow{AM} = \lambda \overrightarrow{AB}\).

\begin{enumerate}
\item 证明： \( \lambda = 1 - e^{2} \)；
\item 若 \(\lambda = \frac{3}{4}\)，\(\triangle PF_{1}F_{2}\) 的周长为 \(6\)；写出椭圆 \(C\) 的方程；
\item 确定 \( \lambda \) 的值，使得 \( \triangle PF_{1}F_{2} \) 是等腰三角形.
\end{enumerate}

\begin{solution}
\begin{enumerate}
\item 记椭圆的焦半距为 \(c\). 由离心率定义，\(c=ae\)，且
\[
b^2=a^2-c^2=a^2(1-e^2),\qquad 0<e<1.
\]
直线 \(l\) 与坐标轴的交点为
\[
A\left(-\frac{a}{e},0\right),\qquad B(0,a).
\]
由 \(\overrightarrow{AM}=\lambda\overrightarrow{AB}\)，得
\[
M\left(\frac{a(\lambda-1)}{e},a\lambda\right).
\]
将点 \(M\) 代入椭圆方程，得
\[
\frac{(\lambda-1)^2}{e^2}+\frac{\lambda^2}{1-e^2}=1.
\]
两边乘以 \(e^2(1-e^2)\)，整理得
\[
(1-e^2)(\lambda-1)^2+e^2\lambda^2=e^2(1-e^2),
\]
即
\[
\left(\lambda-(1-e^2)\right)^2=0.
\]
所以
\[
\lambda=1-e^2.
\]

\item 当 \(\lambda=\frac34\) 时，\(e^2=\frac14\)，由 \(e>0\) 得 \(e=\frac12\). 又
\[
F_1(-ae,0),\qquad F_2(ae,0),
\]
且直线 \(l\) 的方程为 \(ex-y+a=0\). 记
\[
\delta=\frac{e(-ae)-0+a}{1+e^2}=\frac{a(1-e^2)}{1+e^2}.
\]
由点关于直线的对称公式，点 \(F_1\) 的对称点为
\[
P=\left(-ae-2e\delta,2\delta\right)
=\left(-\frac{ae(3-e^2)}{1+e^2},\frac{2a(1-e^2)}{1+e^2}\right).
\]
因此
\[
PF_1=2\operatorname{dist}(F_1,l)=\frac{2a(1-e^2)}{\sqrt{1+e^2}},
\]
\[
PF_2=\frac{2a}{1+e^2}\sqrt{4e^2+(1-e^2)^2}=2a,
\qquad F_1F_2=2ae.
\]
三角形 \(PF_1F_2\) 的周长为
\[
2a\left(\frac{1-e^2}{\sqrt{1+e^2}}+1+e\right).
\]
代入 \(e=\frac12\) 和周长为 \(6\)，得
\[
\frac{3a}{\sqrt5}+2a+a=6,
\]
所以
\[
a=\frac{2\sqrt5}{\sqrt5+1}=\frac{5-\sqrt5}{2}.
\]
又
\[
a^2=\frac{5(3-\sqrt5)}{2},\qquad
b^2=a^2(1-e^2)=\frac{15(3-\sqrt5)}{8}.
\]
故椭圆 \(C\) 的方程为
\[
\frac{x^2}{\frac{5(3-\sqrt5)}{2}}+\frac{y^2}{\frac{15(3-\sqrt5)}{8}}=1.
\]
\item 因为 \(0<e<1\)，所以 \(PF_1<PF_2\)，且 \(PF_2=2a>2ae=F_1F_2\). 因而等腰三角形只能满足 \(PF_1=F_1F_2\)，即
\[
\frac{1-e^2}{\sqrt{1+e^2}}=e.
\]
两边均为正，平方得
\[
(1-e^2)^2=e^2(1+e^2),
\]
即 \(1-3e^2=0\). 由 \(e>0\)，得 \(e=\frac{\sqrt3}{3}\)，所以
\[
\lambda=1-e^2=\frac23.
\]
\end{enumerate}
\end{solution}
\end{problem}
